\setcounter{page}{19}
\chapter*{Chapter I. The Derived Category}

\section*{Introduction}
Let \(\cat{A}\) and \(\cat{B}\) be abelian categories, and let \(F\colon \cat{A}\to\cat{B}\) be a functor. Our purpose is to define, functorially for each \(X\in\operatorname{Ob}\cat{A}\), a complex \(\mathbf{R}F(X)\), whose cohomology groups are the right derived functors of \(F\) acting on \(X\), namely \(R^i F(X)\). In general, for any complex \(X^\bullet\) of objects of \(\cat{A}\), we will define \(\mathbf{R}F(X^\bullet)\).

More precisely, \(\mathbf{R}F\) will be a functor from the derived category \(D(\cat{A})\) of \(\cat{A}\), to the derived category \(D(\cat{B})\) of \(\cat{B}\).

The derived category \(D(\cat{A})\) is obtained as follows: one first considers the category \(K(\cat{A})\), whose objects are complexes of elements of \(\cat{A}\), and whose morphisms are homotopy equivalence classes of morphisms of complexes. The category \(D(\cat{A})\) is obtained by ``localizing'' \(K(\cat{A})\), so that every morphism in \(K(\cat{A})\) which induces an isomorphism on cohomology, becomes an isomorphism in \(D(\cat{A})\). This process of localization will be explained in general below.

Although \(\cat{A}\) and \(\cat{B}\) above are assumed to be abelian categories, the categories \(K(\cat{A})\), \(D(\cat{A})\), \dots are in general not abelian. However, they can be given a structure which carries enough information for our purposes, namely a structure of triangulated category. Thus we will be led to study triangulated categories.